\documentclass[11pt,a4paper]{article}
\usepackage[T1]{fontenc}
\usepackage[utf8]{inputenc}
\usepackage{amsmath,amssymb,mathtools}
\usepackage{booktabs}
\usepackage[hidelinks]{hyperref}
\usepackage[a4paper,margin=2.7cm]{geometry}

\newcommand{\R}{\mathbb{R}}
\newcommand{\grad}{\nabla}
\newcommand{\trans}{\mathsf{T}}
\newcommand{\norm}[1]{\left\lVert #1 \right\rVert}

\title{Line-Search Methods for Optimization\\
\large Theory and relation to the \texttt{LineSearch} example}
\author{PACS examples}
\date{\today}

\begin{document}
\maketitle

\begin{abstract}
This note presents the theory behind the line-search methods implemented in the
\texttt{LineSearch} folder. The code combines interchangeable search-direction
strategies with Armijo backtracking and illustrates a simple projection approach
for component-wise box constraints.
\end{abstract}

\section{Unconstrained minimization and line search}

Consider the smooth optimization problem
\begin{equation}
  \min_{x\in\R^n} f(x).
  \label{eq:problem}
\end{equation}
At an unconstrained local minimizer $x_\star$, a necessary first-order
condition is $\grad f(x_\star)=0$. Line-search methods produce the sequence
\begin{equation}
  x_{k+1}=x_k+\alpha_kd_k,
  \label{eq:update}
\end{equation}
where $d_k$ is a direction and $\alpha_k>0$ is a step length. A direction is
strictly descending if
\begin{equation}
  \grad f(x_k)^\trans d_k<0.
  \label{eq:descent}
\end{equation}
Since
\[
  f(x_k+\alpha d_k)=f(x_k)+\alpha\grad f(x_k)^\trans d_k+o(\alpha),
\]
a sufficiently small positive step decreases the objective. The solver checks
this property before backtracking and falls back to the negative gradient if a
strategy returns a non-descent direction.

The direction and the step length have different roles. The former carries
local first- or second-order information, while the latter controls global
progress and prevents an overly long step from increasing the objective.

\section{Armijo backtracking}

The source file \texttt{LineSearchSolver.cpp} implements geometric
backtracking using the Armijo sufficient-decrease condition:
\begin{equation}
 f(x_k+\alpha d_k)\leq f(x_k)+c_1\alpha\grad f(x_k)^\trans d_k,
 \qquad 0<c_1<1.
 \label{eq:armijo}
\end{equation}
It starts from a trial value $\alpha$, accepts it if \eqref{eq:armijo} holds,
and otherwise replaces it by $\rho\alpha$, with $0<\rho<1$, until acceptance or
a maximum number of reductions.

The corresponding configuration fields are
\texttt{initialStep}, \texttt{sufficientDecreaseCoefficient},
\texttt{stepSizeDecrementFactor}, and \texttt{maxIter} in
\texttt{LineSearchOptions}. Their defaults are $1$, $10^{-2}$, $0.5$, and
$40$. The implementation currently enforces only the first Wolfe condition,
\eqref{eq:armijo}. The option \texttt{secondWolfConditionFactor} is reserved
for the curvature condition
\[
 \left\lvert\grad f(x_k+\alpha d_k)^\trans d_k\right\rvert
 \leq c_2\left\lvert\grad f(x_k)^\trans d_k\right\rvert,
 \qquad c_1<c_2<1,
\]
but that condition is not yet implemented. Strong-Wolfe searches are important
in standard global-convergence theory for BFGS and nonlinear conjugate-gradient
methods.

\section{Directions in the library}

\texttt{DescentDirectionBase} is the common interface for all direction
strategies. Its call operator receives the current point, cost, gradient,
Hessian, and bound data. \texttt{DescentDirectionFactory} maps the names in
Table~\ref{tab:directions} to concrete strategies, enabling selection from
\texttt{linesearch\_options.json}.

\begin{table}[htbp]
\centering
\begin{tabular}{@{}lll@{}}
\toprule
Name & Direction & Information required \\
\midrule
\texttt{GradientDirection} & $-g_k$ & gradient \\
\texttt{NewtonDirection} & $H_kd_k=-g_k$ & gradient and Hessian \\
\texttt{BFGSDirection} & $B_kd_k=-g_k$ & gradient \\
\texttt{BFGSIDirection} & $d_k=-C_kg_k$ & gradient \\
\texttt{BBDirection} & $d_k=-\gamma_kg_k$ & gradient \\
\texttt{CGDirection} & Polak--Ribiere recurrence & gradient \\
\bottomrule
\end{tabular}
\caption{Implemented strategies. Here $g_k=\grad f(x_k)$,
$H_k=\grad^2f(x_k)$, $B_k$ approximates the Hessian, and $C_k$ its inverse.}
\label{tab:directions}
\end{table}

Steepest descent takes $d_k=-g_k$, for which
$g_k^\trans d_k=-\norm{g_k}^2$. It is robust but can be slow on
ill-conditioned problems. Newton's method minimizes the local quadratic model
\[
 f(x_k+p)\approx f(x_k)+g_k^\trans p+\tfrac12p^\trans H_kp
\]
and solves $H_kp=-g_k$. It has local quadratic convergence near a minimizer
with a positive-definite Hessian. The unbounded implementation uses an LLT
solve, so it assumes such a Hessian; production solvers normally safeguard
indefinite Hessians with modifications or trust regions.

For BFGS, define $s_k=x_{k+1}-x_k$ and $y_k=g_{k+1}-g_k$. The direct update is
\begin{equation}
 B_{k+1}=B_k+\frac{y_ky_k^\trans}{y_k^\trans s_k}
 -\frac{B_ks_ks_k^\trans B_k}{s_k^\trans B_ks_k}.
 \label{eq:bfgs}
\end{equation}
When $B_k$ is positive definite and $y_k^\trans s_k>0$, it preserves positive
definiteness. \texttt{BFGSDirection} stores $B_k$, while
\texttt{BFGSIDirection} stores an inverse approximation $C_k$. Both start from
the identity and skip updates that fail a positive-curvature test.

\texttt{BBDirection} uses a scaled gradient based on the two
Barzilai--Borwein estimates
\[
 \gamma_k^{(1)}=\frac{s_k^\trans y_k}{y_k^\trans y_k},\qquad
 \gamma_k^{(2)}=\frac{s_k^\trans s_k}{s_k^\trans y_k}.
\]
It averages them when safe and otherwise reverts to steepest descent.
\texttt{CGDirection} uses the Polak--Ribiere recurrence
\[
 d_{k+1}=-g_{k+1}+\beta_kd_k,\qquad
 \beta_k=\frac{g_{k+1}^\trans(g_{k+1}-g_k)}{g_k^\trans g_k},
\]
and restarts to the negative gradient if this ceases to be descending. BFGS,
BB, and CG retain previous-iteration data; their \texttt{reset()} method is
therefore called when the solver receives a new initial point.

\section{Derivative evaluation}

\texttt{OptimizationData} stores callable cost, gradient, and Hessian
functions. Analytic derivatives can be supplied directly.
\texttt{GradientFiniteDifference} offers forward, backward, and centered
finite-difference gradients. \texttt{GradientsAD} wraps the bundled
\texttt{autodiff} library: a scalar cost written with
\texttt{autodiff::var} variables is converted into ordinary cost, gradient,
and Hessian callables. Reverse mode is well suited to gradients of scalar
objectives with many inputs. The solver may request quantities separately;
\texttt{GradientsAD::evaluateAll()} is available when all are needed together.

\section{Box constraints and projection}

The code can constrain iterates to
\[
 \mathcal{B}=\{x\in\R^n:l_i\leq x_i\leq u_i,\quad i=1,\ldots,n\}.
\]
After \texttt{setBounds()} is called, the initial point must be feasible.
Backtracking projects each trial point component-wise:
\begin{equation}
 x_{k+1}=\Pi_{\mathcal B}(x_k+\alpha_kd_k),\qquad
 [\Pi_{\mathcal B}(z)]_i=\min\{\max\{z_i,l_i\},u_i\}.
 \label{eq:projection}
\end{equation}
This ensures feasible accepted points. The bounded Newton direction also
identifies active components and removes them from its inverse-Hessian
calculation.

For a box-constrained minimizer, ordinary stationarity is replaced by
\begin{equation}
 \grad f(x_\star)^\trans(y-x_\star)\geq0
 \quad\text{for all }y\in\mathcal B.
 \label{eq:box-stationarity}
\end{equation}
Equivalently, the projected gradient vanishes. The present example still uses
the ordinary gradient norm among its stopping tests, and its non-Newton
strategies do not explicitly form projected directions. It is consequently a
compact projection example, not a complete active-set implementation. A robust
extension should add a projected-gradient stopping test, tolerance-aware active
sets, feasible directions for every strategy, or a dedicated method such as
L-BFGS-B.

\section{Reading the implementation}

\begin{itemize}
\item \texttt{LineSearchSolver.hpp/.cpp}: iteration, Armijo backtracking,
      stopping tests, and projection.
\item \texttt{DescentDirectionBase.hpp} and
      \texttt{DescentDirections.hpp/.cpp}: strategy interface and formulas.
\item \texttt{LineSearch\_traits.hpp}, \texttt{Optimization\_options.hpp}, and
      \texttt{LineSearch\_options.hpp}: common types, problem data, and
      configuration.
\item \texttt{main\_linesearch.cpp}: a worked driver.
\end{itemize}

\begin{thebibliography}{9}
\bibitem{NocedalWright}
J.~Nocedal and S.~J.~Wright, \emph{Numerical Optimization}, 2nd ed.,
Springer, 2006.

\bibitem{DennisSchnabel}
J.~E.~Dennis, Jr. and R.~B.~Schnabel,
\emph{Numerical Methods for Unconstrained Optimization and Nonlinear
Equations}, SIAM, 1996.

\bibitem{Fletcher}
R.~Fletcher, \emph{Practical Methods of Optimization}, 2nd ed., Wiley, 1987.

\bibitem{ByrdLuNocedalZhu}
R.~H.~Byrd, P.~Lu, J.~Nocedal, and C.~Zhu,
A limited memory algorithm for bound constrained optimization,
\emph{SIAM Journal on Scientific Computing}, 16(5), pp.~1190--1208, 1995.
\end{thebibliography}

\end{document}
